% paper_18_exchange_constants.tex — GeoVac Paper 18
% Spectral-Geometric Exchange Constants: The Transcendental
% Toll of Projecting Discrete Spectra onto Continuous Manifolds
% Status: Draft, March 29, 2026
% Category: Observations (papers/observations/)
\documentclass[aps,pra,twocolumn,superscriptaddress,longbibliography]{revtex4-2}

\usepackage{amsmath,amssymb,amsthm,graphicx,xcolor,bm,braket,dcolumn,booktabs}
\usepackage{hyperref}

\newtheorem{theorem}{Theorem}
\newtheorem{observation}{Observation}
\newtheorem{conjecture}{Conjecture}

\begin{document}

\title{Spectral-Geometric Exchange Constants:\\
The Transcendental Toll of Projecting Discrete Spectra\\
onto Continuous Manifolds}
\thanks{Paper 18 in the Geometric Vacuum series}

\author{J.~Loutey}
\affiliation{Independent Researcher, Kent, Washington}

\date{March 29, 2026}

%% ========================================================================
%% ABSTRACT
%% ========================================================================

\begin{abstract}
The GeoVac program has identified a recurring pattern across its
hierarchy of natural geometries: each time a discrete algebraic
structure (graph eigenvalues, representation labels, recurrence
relations) is projected onto or embedded in a continuous coordinate
system, a characteristic transcendental constant appears as the
conversion factor.  At Level~1, the projection from $S^3$ to
$\mathbb{R}^3$ introduces the kinetic scale factor
$\kappa = -1/16$.  At Level~2, the embedding of the Laguerre
spectral basis into prolate spheroidal coordinates introduces the
Stieltjes seed $e^a E_1(a)$.  At Level~3, the adiabatic
reparameterization of the $S^3 \times S^3$ product space into
hyperspherical coordinates introduces the transcendental function
$\mu(R)$.  We observe that this pattern is an instance of a
well-known phenomenon in spectral geometry: Weyl's law, which
relates the discrete eigenvalue counting function of the
Laplace--Beltrami operator to the continuous volume of the
underlying manifold, with conversion factors involving
$\pi^{d/2}/\Gamma(d/2+1)$.  The Selberg trace formula generalizes
this to an exact identity between spectral sums and geometric
sums over geodesics, with transcendental conversion factors
determined by the manifold's curvature.  We propose that the
GeoVac constants $\kappa$, $e^a E_1(a)$, and $\mu(R)$ are
specific instances of Weyl--Selberg exchange constants for the
natural geometries $S^3$, prolate spheroid, and $S^5$
respectively.  This identification explains why the constants
are transcendental, why their complexity increases with the
geometry, and why they vanish when the computation remains
on the graph.  We note that the combination rule
$K = \pi(B + F - \Delta)$ in the conjectural fine structure
constant formula (Paper~2) has the same structural character:
$\pi$ multiplying a sum of discrete representation-theoretic
data and continuous spectral invariants.
\end{abstract}

\maketitle

%% ========================================================================
%% I. INTRODUCTION
%% ========================================================================

\section{Introduction}
\label{sec:intro}

The Geometric Vacuum program constructs discrete lattice
Hamiltonians by identifying the natural geometry of each quantum
system and discretizing it as a sparse
graph~\cite{loutey_paper0,loutey_paper1,loutey_paper7}.  A
persistent observation across the hierarchy of natural geometries
is that the discrete algebraic structures---integer eigenvalues,
finite recurrence relations, Gaunt integrals built from Wigner
$3j$ symbols---are entirely self-contained, but that their
connection to physical observables in continuous coordinate
systems always introduces transcendental constants.

This paper catalogs these constants, identifies the mathematical
phenomenon they instantiate, and proposes a unifying
interpretation.

The central observation is simple: each constant is the
\textit{exchange rate} between a discrete spectrum and a
continuous geometry.  The exchange rate is necessarily
transcendental because it encodes the curvature mismatch between
the graph's intrinsic combinatorial structure and the coordinate
system used for physical interpretation.  This observation is
not new in mathematics---it is the content of Weyl's
law~\cite{weyl1911} and the Selberg trace
formula~\cite{selberg1956}---but its systematic appearance in
the GeoVac hierarchy has not been noted, and the connection
to the conjectural fine structure constant formula (Paper~2)
has not been made.


%% ========================================================================
%% II. CATALOG OF EXCHANGE CONSTANTS
%% ========================================================================

\section{Catalog of Exchange Constants}
\label{sec:catalog}

\subsection{Level 1: $\kappa = -1/16$ ($S^3 \to \mathbb{R}^3$)}
\label{sec:kappa}

The graph Laplacian on $S^3$ produces pure integer eigenvalues
\begin{equation}
  \lambda_n = -(n^2 - 1), \qquad n = 1, 2, 3, \ldots
  \label{eq:s3_eigenvalues}
\end{equation}
with degeneracy $g_n = n^2$~\cite{loutey_paper7}.  The physical
hydrogen energy levels $E_n = -1/(2n^2)$ are related by the
energy-shell constraint $p_0^2 = -2E$, which acts as a
stereographic focal length projecting $S^3$ onto flat
$\mathbb{R}^3$.  The universal kinetic scale factor
\begin{equation}
  \kappa = -\frac{1}{16}
  \label{eq:kappa}
\end{equation}
converts graph eigenvalues to physical energies:
$E_n = \kappa \cdot \lambda_n / n^2$~\cite{loutey_paper0}.

The origin of $\kappa$ is well understood geometrically: it is
the Jacobian of the stereographic projection composed with the
conformal factor of the Fock map.  Specifically, if $\Omega$
is the conformal factor $\Omega = (1 + p^2/p_0^2)^{-1}$, then
the Laplace--Beltrami operator on $S^3$ and the flat-space
kinetic operator are related by
$\nabla^2_{S^3} = \Omega^{-2}[\nabla^2_{\mathrm{flat}} + \text{lower order}]$,
and $\kappa$ absorbs the accumulated conformal rescaling and
the $p_0$ energy-shell normalization~\cite{loutey_paper7}.

The key property of $\kappa$ is that it is \textit{unnecessary
on the graph}.  The graph Laplacian eigenvalue problem $L\psi =
\lambda\psi$ is solved entirely in terms of integers.  The
constant $\kappa$ appears only when the result is translated
into atomic units (Hartrees, bohrs)---i.e., when the discrete
spectrum is projected onto a continuous coordinate system.

\subsection{Level 2: $e^a E_1(a)$ (Laguerre basis $\to$ prolate spheroid)}
\label{sec:stieltjes}

The prolate spheroidal radial solver uses a Laguerre spectral
basis $\varphi_n(\xi) = e^{-\alpha(\xi-1)} L_n(2\alpha(\xi-1))$
with matrix elements computed from three-term
recurrences~\cite{loutey_paper11}.  For $\sigma$ states ($m = 0$),
all matrix elements are purely algebraic: the Laguerre moment
matrices $M_k[i,j] = \int_0^\infty x^k L_i(x) L_j(x) e^{-x}
dx$ are band-structured matrices with rational entries, built
from the recurrence $x L_n = -(n+1)L_{n+1} + (2n+1)L_n -
nL_{n-1}$.

For $m \neq 0$ states, the centrifugal term $m^2/(\xi^2 - 1)$
introduces a $1/x$ singularity in the Laguerre variable
$x = 2\alpha(\xi - 1)$.  The associated Laguerre basis
$L_n^{|m|}(x)$ with weight $x^{|m|} e^{-x}$ absorbs this
singularity by construction.  The centrifugal matrix elements
decompose via partial-fraction expansion:
\begin{equation}
  \frac{m^2}{\xi^2 - 1}
  = 4\alpha^2 m^2 \left[\frac{1}{x}
    - \frac{1}{x + 4\beta}\right],
  \label{eq:partial_fraction}
\end{equation}
where $\beta = \alpha(\xi_0 - 1)$ is the Laguerre scale offset.
The two terms have fundamentally different algebraic character:
\begin{itemize}
  \item The $1/x$ term produces a lowered moment
    $M_{-1}[i,j] = \int_0^\infty x^{|m|-1} L_i^{|m|}(x)
    L_j^{|m|}(x)\,e^{-x}\,dx$, which is fully algebraic via
    the DLMF 18.9.13 summation
    identity~\cite{dlmf_laguerre}.
  \item The $1/(x + 4\beta)$ term produces a Stieltjes matrix
    $J[i,j] = \int_0^\infty \frac{x^{|m|}}{x + a}\,
    L_i^{|m|}(x) L_j^{|m|}(x)\,e^{-x}\,dx$, where $a = 4\beta$.
    This matrix satisfies a three-term recurrence in the
    row/column indices, seeded by the single transcendental
\end{itemize}
\begin{equation}
  S_0^{(\alpha)}(a) = \Gamma(\alpha+1)\,e^a\,E_{\alpha+1}(a),
  \label{eq:stieltjes_seed}
\end{equation}
where $E_n(a) = \int_1^\infty t^{-n} e^{-at}\,dt$ is the
generalized exponential integral and $a = 4\beta$ is the
Laguerre scale parameter~\cite{loutey_track_j}.  For integer
$\alpha = |m|$, all higher $E_n$ reduce to $E_1(a)$ via a
finite rational recurrence, so the entire matrix is built from
algebraic operations plus one transcendental constant.

The computational cost of the transcendental seed is negligible:
one call to $E_1(a)$, which achieves machine precision for
$a \ge 4$ (covering all physical parameters).  The seed then
propagates through an $O(N^2)$ algebraic recurrence that produces
all matrix elements.  The continuous mathematics has been
compressed from a \textit{function} (the quadrature integrand
over $[0,\infty)$) to a \textit{constant} (the Stieltjes seed),
with the entire function-level information reconstructed
algebraically from that constant.

An unexpected computational bonus: the associated Laguerre basis
converges \textit{faster} than the ordinary Laguerre basis for
$m \neq 0$ states.  For $|m| = 1$ ($\pi$ states), the associated
basis is stable by $N = 10$, compared to non-monotonic convergence
at $N \approx 20$ for the ordinary basis.  For $|m| = 2$
($\delta$ states), algebraic-quadrature agreement reaches
$\sim 5 \times 10^{-9}$ (machine precision), with PES shape
(equilibrium geometry) preserved exactly~\cite{loutey_track_j}.
This suggests that identifying the exchange constant and absorbing
it into the basis does not merely explain the transcendence---it
\textit{improves the spectral method}, because the adapted basis
eliminates the curvature mismatch that the exchange constant
mediates.

The Stieltjes seed $e^a E_1(a)$ is the Laplace transform of the
rational function $1/(1+t)$ evaluated at $a$.  It encodes how
the Laguerre inner product (defined on the flat half-line
$[0, \infty)$ with exponential weight) ``sees'' a shifted pole
at $x = -a$ arising from the curvature of the prolate spheroidal
coordinate system.  On the graph, there is no pole; the
centrifugal barrier is a selection rule encoded in the quantum
numbers.  The transcendental appears only when this selection
rule is represented as a rational function in a continuous
coordinate.

\subsection{Level 3: $\mu(R)$ ($S^3 \times S^3 \to$ hyperspherical)}
\label{sec:mu}

The two-electron problem on $S^3 \times S^3$ is fully algebraic:
the Hamiltonian matrix has entries given by Gaunt integrals
(finite sums of Wigner $3j$ symbols with rational coefficients),
and the eigenvalues are roots of a polynomial with algebraic
coefficients.  The Full CI solver exploits this, achieving
0.35\% error for helium with no transcendentals
anywhere~\cite{loutey_fci}.

The hyperspherical reparameterization introduces coordinates
$(R, \alpha)$ where $R = \sqrt{r_1^2 + r_2^2}$ is the
hyperradius and $\alpha = \arctan(r_2/r_1)$ is the hyperangle.
The angular eigenvalues $\mu(R)$ depend parametrically on $R$
because the effective angular Hamiltonian is
$H(R) = H_0 + R \cdot V^{\mathrm{coupling}}$~\cite{loutey_paper13}.

Each perturbation coefficient $a_k$ in the expansion
$\mu(R) = \mu^{(0)} + a_1 R + a_2 R^2 + \cdots$ is individually
algebraic: the $a_k$ are finite sums of rational numbers times
powers of $\pi^{-1}$, arising from partial harmonic sums over
$\mathrm{SO}(6)$ Casimir eigenstates~\cite{loutey_paper13}.
However, the full function $\mu(R)$ is transcendental: it
interpolates between $O(R)$ behavior at small $R$ (perturbative,
$S^5$ symmetric) and $O(R^2)$ behavior at large $R$
(hydrogen-like, $S^3 \times \mathbb{R}$), and no finite-order
rational function spans both
regimes~\cite{loutey_paper13,loutey_track_g}.

The transcendence of $\mu(R)$ reflects a topology change in
the underlying manifold: at $R \to 0$, the angular manifold
is the round five-sphere $S^5$ with $\mathrm{SO}(6)$ symmetry;
at $R \to \infty$, it pinches to $S^3 \times \mathbb{R}$
(one electron bound, one free)~\cite{loutey_paper13}.  This
smooth deformation of the manifold as a function of the
``slow'' variable $R$ produces a transcendental function of
$R$ because smooth topology changes are generically
transcendental functions of the deformation parameter.

The crucial point is that $\mu(R)$ is not intrinsic to the
two-electron problem.  It is an artifact of the hyperspherical
coordinate system.  The $S^3 \times S^3$ product space does not
require $\mu(R)$; it only appears when one reparameterizes for
computational efficiency.  The exchange rate between the discrete
product-space description and the continuous hyperspherical
description is an entire transcendental function, not merely
a constant, because the hyperspherical embedding distorts the
product geometry differently at each $R$.


%% ========================================================================
%% III. THE PATTERN
%% ========================================================================

\section{The Pattern: Weyl--Selberg Exchange}
\label{sec:pattern}

\subsection{Weyl's law}

The mathematical content of the observations in
Sec.~\ref{sec:catalog} is an instance of Weyl's law, which
states that the eigenvalue counting function $N(\lambda)$ of
the Laplace--Beltrami operator on a compact Riemannian manifold
of dimension $d$ satisfies~\cite{weyl1911}
\begin{equation}
  N(\lambda) \sim \frac{\omega_d}{(2\pi)^d}\,
  \mathrm{vol}(\Omega)\,\lambda^{d/2}
  \qquad (\lambda \to \infty),
  \label{eq:weyl}
\end{equation}
where $\omega_d = \pi^{d/2}/\Gamma(d/2+1)$ is the volume of
the unit ball in $\mathbb{R}^d$.

Weyl's law relates two fundamentally different descriptions
of the same manifold:
\begin{itemize}
  \item The \textbf{spectral side}: $N(\lambda)$, a discrete
    staircase function counting integer-labeled eigenvalues.
  \item The \textbf{geometric side}: $\mathrm{vol}(\Omega)$, a
    continuous quantity measured by integration over the manifold.
\end{itemize}
The conversion factor $\omega_d / (2\pi)^d$ involves
$\pi^{d/2}$---a transcendental that depends only on the
dimension.  This factor is the prototypical
\textit{spectral-geometric exchange constant}: it converts
between counting (discrete) and measuring (continuous).

For the specific manifolds in the GeoVac hierarchy:
\begin{itemize}
  \item $S^1$ ($d = 1$): $\omega_1 = 2$,
    exchange factor $\sim \pi^{1/2}$.
  \item $S^2$ ($d = 2$): $\omega_2 = \pi$,
    exchange factor $\sim \pi$.
  \item $S^3$ ($d = 3$): $\omega_3 = 4\pi/3$,
    exchange factor $\sim \pi^{3/2}$.
  \item $S^5$ ($d = 5$): $\omega_5 = 8\pi^2/15$,
    exchange factor $\sim \pi^{5/2}$.
\end{itemize}

\subsection{$\pi$ as the founding example}

The simplest spectral-geometric exchange constant is $\pi$ itself.
The Laplacian on $S^1$ (the circle) has eigenvalues $n^2$, counted
by integers; the continuous circumference is $2\pi$.  The
conversion between the discrete eigenvalue count and the continuous
measure is mediated by $\pi$.  More generally, the volume of the
unit $d$-sphere involves $\pi^{d/2}$, and the Weyl constant
$\omega_d / (2\pi)^d$ carries exactly this factor.  Every exchange
constant in the GeoVac hierarchy is a descendant of this primordial
example: $\kappa$ inherits from the $S^3$ Weyl constant, the
Stieltjes seed involves the exponential integral (a function whose
Taylor coefficients involve the Euler--Mascheroni constant, itself
a limit of harmonic sums weighted by $\pi$-dependent terms), and
$\mu(R)$ is an eigenvalue of a Hamiltonian on $S^5$ whose Weyl
asymptotics carry $\pi^{5/2}$.

\subsection{The Selberg trace formula}

Weyl's law is asymptotic.  The Selberg trace
formula~\cite{selberg1956} provides the exact identity:
\begin{equation}
  \sum_j h(r_j)
  = \frac{\mathrm{Area}(\Gamma \backslash \mathcal{H})}{4\pi}
    \int_{-\infty}^{\infty} r\,\tanh(\pi r)\,h(r)\,dr
    + \sum_{\{\gamma\}} \cdots
  \label{eq:selberg}
\end{equation}
where the left side sums over the discrete spectrum, the first
term on the right involves the continuous volume (with a factor
of $\pi$), and subsequent terms sum over conjugacy classes
(geodesics).  The Selberg trace formula is the bridge between
the spectral world and the geometric world, and the
transcendental factors ($\pi$, $\tanh$, logarithms of
algebraic numbers) are the toll for crossing.

When the quotient $\Gamma \backslash G$ is compact, the spectrum
is purely discrete and the trace formula involves only algebraic
data plus $\pi$-dependent conversion factors.  When the quotient
is non-compact, a continuous spectrum appears and the exchange
involves Eisenstein series---transcendental
functions~\cite{selberg1956}.  This exactly parallels the
GeoVac hierarchy: Level~1 (compact $S^3$) has a purely algebraic
spectrum plus a constant ($\kappa$); Level~3 (non-compact
$S^5 \to S^3 \times \mathbb{R}$ flow) has a transcendental
function ($\mu(R)$).


%% ========================================================================
%% IV. TAXONOMY
%% ========================================================================

\section{Taxonomy of Exchange Constants}
\label{sec:taxonomy}

Based on the catalog and the Weyl--Selberg identification, we
propose the following taxonomy:

\begin{table}[h]
\caption{Spectral-geometric exchange constants in the GeoVac
  hierarchy.  ``Source'' indicates the discrete structure;
  ``Target'' indicates the continuous coordinate system.
  ``Type'' classifies the exchange constant.}
\label{tab:taxonomy}
\begin{ruledtabular}
\begin{tabular}{ccccc}
  Level & Source & Target & Constant & Type \\
  \hline
  --- & $S^1$ lattice & $\mathbb{R}^1$ & $\pi$ & Weyl \\
  1 & $S^3$ graph & $\mathbb{R}^3$ & $\kappa = -1/16$ & conformal \\
  2 ($m{=}0$) & Laguerre rec. & prolate sph. & (none) & algebraic \\
  2 ($m{\neq}0$) & assoc. Laguerre & prolate sph. & $e^a E_1(a)$
    & Stieltjes \\
  3 (FCI) & $S^3 {\times} S^3$ & (none) & (none) & algebraic \\
  3 (hyper.) & $\mathrm{SO}(6)$ reps & $S^5(R)$ flow
    & $\mu(R)$ & Eisenstein \\
\end{tabular}
\end{ruledtabular}
\end{table}

The taxonomy reveals a hierarchy of exchange types:
\begin{enumerate}
  \item \textbf{Algebraic} (no exchange constant): When the
    computation stays on the graph or uses a basis whose inner
    product matches the graph structure, no transcendental
    appears.  Examples: Level~1 on $S^3$, Level~2 $\sigma$
    states, Level~3 FCI.

  \item \textbf{Weyl constants} ($\pi^{d/2}$): The
    volume-to-eigenvalue conversion factors for compact
    manifolds.  These are universal---they depend only on
    dimension, not on the specific manifold.  Example: $\pi$
    itself, for $S^1 \to \mathbb{R}$.

  \item \textbf{Conformal constants}: Specific to the conformal
    map between two manifolds of different curvature.  Example:
    $\kappa = -1/16$ for the Fock stereographic projection
    $S^3 \to \mathbb{R}^3$.

  \item \textbf{Stieltjes constants}: Single transcendental
    numbers that seed algebraic recurrences.  These appear when
    a rational singularity in the coordinate system is not
    representable in the spectral basis.  Example:
    $e^a E_1(a)$ for the prolate centrifugal pole.

  \item \textbf{Eisenstein functions}: Transcendental functions
    of a continuous parameter, arising when the manifold itself
    varies.  These appear when the geometry flows (topology
    changes, avoided crossings, non-compact limits).  Example:
    $\mu(R)$ for the $S^5 \to S^3 \times \mathbb{R}$ topology
    flow.
\end{enumerate}

The complexity of the exchange constant increases monotonically
with the severity of the curvature mismatch between source and
target.


%% ========================================================================
%% V. THE FINE STRUCTURE CONSTANT CONNECTION
%% ========================================================================

\section{Connection to the Fine Structure Constant}
\label{sec:alpha}

Paper~2~\cite{loutey_paper2} observes that the fine structure
constant $\alpha$ satisfies the cubic $\alpha^3 - K\alpha + 1 = 0$,
where the coefficient
\begin{equation}
  K = \pi\!\left(42 + \frac{\pi^2}{6} - \frac{1}{40}\right)
  = 137.036064
  \label{eq:K}
\end{equation}
is constructed from spectral invariants of the Hopf fibration
$S^1 \to S^3 \to S^2$.  The combination rule
$K = \pi(B + F - \Delta)$ was identified as conjectural---an
empirical formula with structural support but no first-principles
derivation.

We observe that this combination rule has precisely the structure
of a spectral-geometric exchange:
\begin{itemize}
  \item $B = 42$ is discrete representation-theoretic data:
    the degeneracy-weighted $\mathrm{SO}(3)$ Casimir trace over
    the first three shells of $S^3$.
  \item $F = \zeta(2) = \pi^2/6$ is a continuous spectral
    invariant: the spectral zeta function of $S^1$ at $s = 1$.
  \item $\Delta = 1/40$ is a boundary correction mixing discrete
    eigenvalue data ($|\lambda_3| = 8$) with the cumulative state
    count ($N(2) = 5$).
  \item The overall factor $\pi$ is the Weyl exchange constant
    for $S^2$ (the base of the Hopf fibration, the photon
    propagation sphere).
\end{itemize}

The structure of $K$ is thus: the $S^2$ Weyl constant ($\pi$)
multiplying a sum of the $S^2$ Casimir trace ($B$, discrete),
the $S^1$ spectral zeta value ($F$, transcendental), and an
$S^3$ boundary term ($\Delta$, mixed).  This is a Weyl-type
formula that converts between the discrete bundle
(representation labels on $S^3$) and the continuous coupling
(the electromagnetic interaction strength measured in
$\mathbb{R}^3$).

If this identification is correct, then the combination rule
$K = \pi(B + F - \Delta)$ is not arbitrary numerology but an
instance of the same spectral-geometric exchange that produces
$\kappa = -1/16$ at Level~1 and $e^a E_1(a)$ at Level~2.
The exchange constant for the Hopf bundle is $K$ itself---the
``cost'' of projecting the $S^3$ electron state space through
the $S^1$ gauge fiber onto the $S^2$ photon base, measured
in units of coupling strength.

\begin{conjecture}
The combination rule $K = \pi(B + F - \Delta)$ is a
spectral-geometric exchange constant for the Hopf fibration,
analogous to the Weyl--Selberg constants for compact manifolds.
The specific form---$\pi$ multiplying a sum of Casimir data
and zeta values---should be derivable from the spectral geometry
of $S^1 \to S^3 \to S^2$, potentially via the heat kernel
expansion or the analytic torsion of the Hopf bundle.
\end{conjecture}

Two additional observations support this identification:

\begin{observation}[Structural match]
The combination rule $K = \pi(B + F - \Delta)$ mixes discrete
data ($B = 42$, from representation theory) with continuous
spectral data ($F = \zeta(2) = \pi^2/6$, from the $S^1$
spectral zeta function), multiplied by an overall geometric
factor ($\pi$, the Weyl exchange constant for $S^2$, the Hopf
base).  This is precisely the pattern identified throughout this
paper: a transcendental conversion factor mediating between
discrete and continuous descriptions of the same geometry.
The structural match is evidence---not proof---that the
combination rule is an instance of the exchange constant pattern
rather than numerology.
\end{observation}

\begin{observation}[$\alpha$ is likely transcendental]
The fine structure constant $\alpha$ is a root of the cubic
$\alpha^3 - K\alpha + 1 = 0$, where $K$ involves $\pi$ and
$\zeta(2) = \pi^2/6$.  Roots of polynomials with transcendental
coefficients \textit{can} be algebraic (e.g., $x^2 - \pi x +
\pi(\pi - 1) = 0$ has the algebraic root $x = 1$), but this
requires special cancellations that are generically absent.
Since $K$ is a non-trivial combination of $\pi$ and $\pi^2$
with rational coefficients, it is overwhelmingly likely
(though not proven) that $\alpha$ is transcendental---consistent
with its role as a spectral-geometric exchange constant.
\end{observation}


%% ========================================================================
%% VI. DISCUSSION
%% ========================================================================

\section{Discussion}
\label{sec:discussion}

\subsection{Why the constants are transcendental}

The fundamental reason that spectral-geometric exchange constants
are transcendental is that they convert between two
incommensurable descriptions: a discrete count (eigenvalues,
representation labels, graph vertices) and a continuous measure
(volume, curvature, geodesic length).  The Lindemann--Weierstrass
theorem guarantees that $e^\alpha$ is transcendental for any
nonzero algebraic $\alpha$; since the Weyl constants involve
exponentials of curvature data (via the heat kernel
$\mathrm{tr}\,e^{-t\Delta}$), transcendence is
generic~\cite{berger2003}.

\subsection{Why the constants vanish on the graph}

When the computation remains entirely on the discrete graph---no
projection, no embedding, no coordinate change---the exchange
constant is unnecessary.  The Level~1 $S^3$ graph has integer
eigenvalues; $\kappa$ is needed only to convert them to Hartrees.
The Level~3 FCI on $S^3 \times S^3$ has algebraic matrix elements;
$\mu(R)$ appears only in the hyperspherical reparameterization.
The Level~2 $\sigma$-state solver has algebraic Laguerre moments;
$e^a E_1(a)$ appears only for $m \neq 0$ states where the
coordinate curvature introduces a pole.

This suggests a design principle for the GeoVac program:
\textit{minimize the number of coordinate projections} to minimize
the number of exchange constants.  The ideal solver stays on the
graph; every departure from the graph incurs a transcendental
toll.

\subsection{Exchange constants are computationally cheap}

A striking feature of the exchange constants is that their
conceptual complexity does not correlate with computational cost:
\begin{itemize}
  \item $\kappa = -1/16$: one multiplication.
  \item $e^a E_1(a)$: one function evaluation, followed by an
    $O(N^2)$ algebraic recurrence.
  \item $\mu(R)$: a $10$--$30$ dimensional matrix diagonalization
    per $R$-point (milliseconds).
\end{itemize}
The hierarchy of transcendence---constant, seed, function---tracks
the severity of curvature mismatch, not computational difficulty.
This reinforces the interpretation that exchange constants encode
\textit{geometric information} (the shape of the mismatch between
discrete and continuous descriptions) rather than
\textit{computational obstacles}.  The Level~2 result is
particularly instructive: identifying the Stieltjes seed and
absorbing it into the basis actually \textit{reduces} the
computational cost (by eliminating quadrature) while simultaneously
providing better convergence (Sec.~\ref{sec:stieltjes}).

\subsection{The hierarchy of transcendence}

The exchange constants increase in complexity as the curvature
mismatch increases:
\begin{itemize}
  \item Constant $\to$ constant (Level 1): $\kappa = -1/16$.
    One projection, fixed curvature.
  \item Constant $\to$ single transcendental (Level 2, $m \neq 0$):
    $e^a E_1(a)$.  One embedding with a coordinate singularity.
  \item Constant $\to$ transcendental function (Level 3, hyperspherical):
    $\mu(R)$.  A continuously deforming manifold.
\end{itemize}

This hierarchy mirrors the mathematical classification of periods
and transcendental functions: the Weyl constants are periods
($\pi = \int dx/(1+x^2)$); the Stieltjes seeds are exponential
periods ($E_1(a) = \int_1^\infty e^{-at}/t\,dt$); and the
adiabatic eigenvalue functions are solutions of parameter-dependent
spectral problems, which are generically non-period transcendentals.


%% ========================================================================
%% VII. CONCLUSION
%% ========================================================================

\section{Conclusion}
\label{sec:conclusion}

The GeoVac hierarchy of natural geometries produces a systematic
sequence of transcendental constants at each level where a
discrete algebraic structure is projected onto a continuous
coordinate system.  These constants---$\kappa = -1/16$,
$e^a E_1(a)$, and $\mu(R)$---are instances of the
spectral-geometric exchange constants described by Weyl's law
and the Selberg trace formula.  They encode the curvature
mismatch between the graph and the coordinate system, and they
vanish when the computation stays on the graph.

The observation that the conjectural fine structure constant
formula $K = \pi(B + F - \Delta)$ has the same structural
character---a Weyl-type conversion factor mixing discrete
representation data with continuous spectral invariants---suggests
that the combination rule may be derivable from the spectral
geometry of the Hopf bundle, rather than being an irreducible
empirical observation.

Two implications for the GeoVac program follow:
\begin{enumerate}
  \item \textbf{The design principle}: Staying on the graph
    is not merely a computational preference; it is the
    algebraically natural thing to do.  Every coordinate
    projection incurs a transcendental toll.  The program's
    ongoing effort to algebraicize the solvers
    (Paper~12's Neumann expansion, the Laguerre moment
    recurrences, the Gaunt integral replacements) is, in
    retrospect, an effort to eliminate exchange constants by
    eliminating projections.
  \item \textbf{The physical interpretation of $\kappa$}:
    Paper~0 identified $\kappa = -1/16$ as a universal packing
    constant but did not derive it from first principles.  The
    Weyl--Selberg identification suggests that $\kappa$ is the
    conformal exchange constant for the Fock projection
    $S^3 \to \mathbb{R}^3$, computable from the heat kernel
    expansion of the $S^3$ Laplacian.  This is a testable
    prediction.
\end{enumerate}

\begin{acknowledgments}
Computational implementation and documentation drafting were
performed with the assistance of large language models (Anthropic
Claude) under the author's scientific direction.  All results
were validated against analytical benchmarks.
\end{acknowledgments}


%% ========================================================================
%% BIBLIOGRAPHY
%% ========================================================================

\begin{thebibliography}{30}

\bibitem{loutey_paper0}
J.~Loutey,
``The Geometric Atom: Quantum State Space as a Packing Problem,''
GeoVac Paper~0 (2026).

\bibitem{loutey_paper1}
J.~Loutey,
``The Geometric Atom: Quantum Mechanics as a Packing Problem,''
GeoVac Paper~1 (2026).

\bibitem{loutey_paper7}
J.~Loutey,
``The Dimensionless Vacuum: Recovering the Schr\"odinger Equation
from Scale-Invariant Graph Topology,''
GeoVac Paper~7 (2026).

\bibitem{loutey_paper11}
J.~Loutey,
``The Molecular Fock Projection: Diatomic Molecules as Prolate
Spheroidal Lattices,''
GeoVac Paper~11 (2026).

\bibitem{loutey_paper13}
J.~Loutey,
``The Hyperspherical Lattice: Two-Electron Atoms as Coupled
Channel Graphs,''
GeoVac Paper~13 (2026).

\bibitem{loutey_fci}
J.~Loutey,
``Full configuration interaction on a sparse graph Hamiltonian:
Multi-electron atoms via relational lattice indexing,''
GeoVac FCI Paper (2026).

\bibitem{loutey_paper2}
J.~Loutey,
``The Fine Structure Constant from Spectral Geometry of the Hopf
Fibration,''
GeoVac Paper~2 (2026).

\bibitem{loutey_track_j}
J.~Loutey,
``Track J: Algebraic Level 2 Radial Solver for $m \neq 0$ States,''
GeoVac supplementary material (2026).

\bibitem{loutey_track_g}
J.~Loutey,
``Track G: Perturbation Series for $\mu(R)$,''
GeoVac supplementary material (2026).

\bibitem{dlmf_laguerre}
NIST Digital Library of Mathematical Functions,
\S18.9.13, \url{https://dlmf.nist.gov/18.9.13}.

\bibitem{weyl1911}
H.~Weyl,
``\"Uber die asymptotische Verteilung der Eigenwerte,''
Nachr.\ K\"onigl.\ Ges.\ Wiss.\ G\"ottingen, Math.-phys.\ Kl.\
(1911), pp.\ 110--117.

\bibitem{selberg1956}
A.~Selberg,
``Harmonic analysis and discontinuous groups in weakly symmetric
Riemannian spaces with applications to Dirichlet series,''
J.\ Indian Math.\ Soc.\ \textbf{20}, 47--87 (1956).

\bibitem{berger2003}
M.~Berger,
\textit{A Panoramic View of Riemannian Geometry}
(Springer, Berlin, 2003).

\end{thebibliography}

\end{document}
